\section{模型构建}
\subsection{建模目标与总体框架}
本文的建模目标包含两个互补维度：一是解释维度，即识别碳排放变化的主要驱动机制与作用方向；二是预测维度，即在真实外推场景下比较不同模型的稳定性，并据此开展政策情景模拟。围绕该目标，本文采用“EKC检验 + 扩展STIRPAT + 中介分析 + 滚动预测 + 情景模拟”的组合框架，以保证机制识别与路径预测在同一口径下衔接。

该框架的逻辑顺序为：先通过EKC检验判断经济发展与排放关系是否存在阶段性，再利用扩展STIRPAT估计多因素边际影响，随后用中介框架讨论能源结构影响的传导渠道；在解释基础上，再以滚动一步预测比较模型适配性，最后通过情景设定将统计结果转化为可比较的政策路径。

\subsection{环境库兹涅茨曲线检验}
为检验经济增长与碳排放之间的非线性关系，本文构建对数双次项EKC模型：
\begin{equation}
\ln C_t = \alpha_0 + \alpha_1 \ln A_t + \alpha_2 (\ln A_t)^2 + \varepsilon_t,
\end{equation}
其中$C_t$为第$t$年全国碳排放总量（Mt），$A_t$为人均GDP（万元），$\varepsilon_t$为随机误差项。对两侧变量取对数后，$\alpha_1$可解释为人均GDP对碳排放的弹性，$\alpha_2$捕捉该弹性随收入水平的变化方向。

EKC假说的理论内涵可由排放对人均收入的条件边际效应刻画：
\begin{equation}
\frac{\partial \ln C_t}{\partial \ln A_t} = \alpha_1 + 2\alpha_2 \ln A_t.
\end{equation}
当$\alpha_1>0$且$\alpha_2<0$时，上式随收入增长单调递减，意味着存在排放强度”脱钩拐点”，对应人均GDP的峰值点为：
\begin{equation}
A^* = \exp\!\left(-\frac{\alpha_1}{2\alpha_2}\right).
\end{equation}
若$A^*$落在样本范围内，则支持样本期内存在倒U型关系；若两系数同号或$A^*$远超样本上限，则说明当前阶段排放与收入仍呈单调关系。

EKC检验在本文中的作用并非单独给出因果结论，而是判断”经济规模扩张是否已出现边际减排拐点迹象”，为后续机制解释提供阶段性背景。该定位可避免将简单二次拟合过度解读为结构性因果关系。

{\songti\bfseries{}注：受限于样本长度与对数变换后的极值平滑效应，图7采用水平值二次多项式进行可视化趋势拟合，以更直观地展示拐点形态；公式(6)的对数形式用于理论机制说明与弹性计算基准。}

\subsection{扩展STIRPAT模型}
经典IPAT恒等式将环境压力$I$表达为人口$P$、富裕度$A$与技术水平$T$的乘积：
\begin{equation}
I = P \cdot A \cdot T.
\end{equation}
STIRPAT（Stochastic Impacts by Regression on Population, Affluence and Technology）对其进行随机化扩展：
\begin{equation}
I = a P^b A^c T^d \cdot e,
\end{equation}
其中$a$为模型系数，$b,c,d$为各因素的弹性参数，$e$为随机误差。对两侧取自然对数，得到可估计的对数线性形式：
\begin{equation}
\ln I_t = \ln a + b\ln P_t + c\ln A_t + d\ln T_t + \ln e_t.
\end{equation}
此时$b,c,d$分别表示人口、富裕度和技术对环境压力的弹性，便于直接比较各驱动因素的相对贡献。

结合我国碳排放的现实机制，本文在标准三因素框架基础上引入能源结构、城镇化和产业结构，扩展为六因素模型：
\begin{equation}
\ln C_t = \beta_0 + \beta_1 \ln A_t + \beta_2 \ln EI_t + \beta_3 CS_t + \beta_4 UR_t + \beta_5 TP_t + \beta_6 IS_t + \varepsilon_t,
\end{equation}
各变量含义及预期符号如表~\ref{tab:stirpat_vars}所示。

\begin{table}[H]
\centering
\caption{扩展STIRPAT模型变量说明}
\label{tab:stirpat_vars}
\begin{tabular}{cccc}
\toprule
符号 & 含义 & 单位/度量 & 预期符号 \\
\midrule
$C_t$  & 全国碳排放总量    & Mt         & —（因变量）\\
$A_t$  & 人均GDP（富裕度） & 万元/人    & $+$ \\
$EI_t$ & 能源强度          & 万吨标煤/亿元 & $+$ \\
$CS_t$ & 煤炭消费占比      & 比例        & $+$ \\
$UR_t$ & 城镇化率          & 比例        & $\pm$ \\
$TP_t$ & 技术进步代理（有效发明专利对数） & — & $-$ \\
$IS_t$ & 第二产业增加值占比 & 比例       & $+$ \\
\bottomrule
\end{tabular}
{\footnotesize 注：技术进步预期符号为负，代表其具备减排潜力，但在小样本阶段可能受能源反弹效应干扰表现不显著。}
\end{table}

参数$\beta_1,\beta_2$的估计值可解释为弹性：例如$\hat\beta_1=0.59$表示人均GDP每增长1\%，碳排放平均增长0.59\%。$\beta_3$至$\beta_6$对应水平变量，其系数反映该变量每变动一单位时碳排放对数的变化量（半弹性）。

模型采用普通最小二乘法（OLS）估计，置信区间基于残差方差和$(X^\top X)^{-1}$矩阵的对角元素构造：
\begin{equation}
\hat\beta_j \pm z_{0.975}\cdot\hat{\sigma}\sqrt{[(X^\top X)^{-1}]_{jj}},
\end{equation}
其中$\hat\sigma^2 = \frac{\mathbf{e}^\top\mathbf{e}}{n-k}$，$n=18$为观测数，$k=7$为参数个数（含截距），自由度$n-k=11$。

该设定具有三点优势：第一，避免将人口规模与人均GDP同时引入导致的恒等式相关性干扰，降低多重共线性风险；第二，弹性系数便于比较不同变量的相对作用强度；第三，与政策变量具有直接映射，便于后续情景参数解释。

\subsection{中介效应分析}
中介效应分析用于检验自变量$X$对因变量$Y$的影响是否通过中间变量$M$部分或完全传导。本文以煤炭占比（$CS_t$）作为核心自变量$X_t$，技术进步（$TP_t$）作为潜在中介变量$M_t$，碳排放（$C_t$）的对数作为因变量$Y_t$，依次估计以下三个方程：

{\songti\bfseries{}第一步：}检验$X$对中介变量$M$的影响（$a$路径）：
\begin{equation}
M_t = \gamma_0 + \gamma_1 X_t + u_t.
\end{equation}
$\gamma_1$显著则说明能源结构对技术进步存在可识别关联。

{\songti\bfseries{}第二步：}检验$X$对$Y$的总效应（$c$路径）：
\begin{equation}
Y_t = \delta_0 + \delta_1 X_t + v_t.
\end{equation}
$\delta_1$为自变量对因变量的总效应系数。

{\songti\bfseries{}第三步：}在控制中介变量后，同时估计$X$对$Y$的直接效应（$c'$路径）与$M$对$Y$的效应（$b$路径）：
\begin{equation}
Y_t = \theta_0 + \theta_1 X_t + \theta_2 M_t + w_t.
\end{equation}

中介效应的分解关系为：
\begin{equation}
\underbrace{\delta_1}_{\text{总效应}} = \underbrace{\theta_1}_{\text{直接效应}} + \underbrace{\gamma_1 \cdot \theta_2}_{\text{间接效应（中介）}}.
\end{equation}
若$\gamma_1$与$\theta_2$均显著，且$|\theta_1| < |\delta_1|$，则可认为$M$在$X\to Y$路径上存在部分中介作用；若$\theta_1$不再显著而$\gamma_1\cdot\theta_2$显著，则为完全中介。

需要强调，中介框架在本文中用于揭示”作用链可能性”而非严格因果识别。受限于样本长度，本文更重视路径方向的一致性和经济含义的可解释性。

\subsection{滚动预测模型设计与评估}
为贴近真实应用场景，本文采用滚动一步向前预测（rolling one-step-ahead forecast）而非全样本内拟合。设共有$T$期观测，在预测第$t$期时，仅使用$\{1,2,\ldots,t-1\}$期数据估计模型参数，再预测$\hat{C}_t$，随后将窗口向前滚动一期。该策略严格避免未来信息泄露，所得误差指标更能反映模型的真实外推能力。

本文比较七类模型，各模型的数学形式与参数含义如下：

\begin{enumerate}

\item {\songti\bfseries{}朴素模型（Naive）}：以上一期实际值作为当期预测值，是时序预测的基准参照：
\begin{equation}
\hat{C}_t = C_{t-1}.
\end{equation}
该模型无需估计任何参数，捕捉的是"随机游走"假设下的零信息增量。

\item {\songti\bfseries{}二阶自回归模型（AR2）}：利用前两期实际值线性外推：
\begin{equation}
\hat{C}_t = \phi_0 + \phi_1 C_{t-1} + \phi_2 C_{t-2},
\end{equation}
其中$\phi_0,\phi_1,\phi_2$通过OLS在训练窗口内估计。AR2隐含假设序列的短期动态可由两阶自相关结构完整描述，适用于具有惯性特征的宏观总量序列。

\item {\songti\bfseries{}二次趋势模型（Quadratic Trend）}：以时间$t$为解释变量，拟合确定性二次趋势：
\begin{equation}
\hat{C}_t = \eta_0 + \eta_1 t + \eta_2 t^2,
\end{equation}
其中$\eta_0,\eta_1,\eta_2$通过训练窗口内的多项式回归估计。$\eta_2>0$对应加速上升，$\eta_2<0$对应增速放缓或潜在拐点。

\item {\songti\bfseries{}多变量回归模型（Multivariate OLS）}：引入宏观驱动变量作为解释项：
\begin{equation}
\hat{C}_t = \lambda_0 + \sum_{k=1}^{K} \lambda_k X_{k,t},
\end{equation}
其中$X_{k,t}$包含GDP总量、能源消费总量、煤炭占比、城镇化率、技术进步代理和第二产业占比共$K=6$个变量，参数$\boldsymbol{\lambda}$通过OLS估计：
\begin{equation}
\hat{\boldsymbol{\lambda}} = (X^\top X)^{-1}X^\top \mathbf{C}.
\end{equation}

\item {\songti\bfseries{}随机森林回归（Random Forest）}：通过对$B$棵独立决策树的预测结果取均值以降低方差：
\begin{equation}
\hat{C}_t^{(\mathrm{RF})} = \frac{1}{B}\sum_{b=1}^{B} T_b(\mathbf{X}_t),
\end{equation}
其中$T_b(\cdot)$为第$b$棵回归树，每棵树在自助抽样子集与随机特征子集上训练。特征随机化减少树间相关，有助于提高集成预测稳定性，但在小样本下易过拟合。

\item {\songti\bfseries{}XGBoost梯度提升树}：通过加法方式逐步拟合残差，最终预测为$M$棵树之和：
\begin{equation}
\hat{C}_t^{(\mathrm{XGB})} = \sum_{m=1}^{M} f_m(\mathbf{X}_t), \quad f_m \in \mathcal{F},
\end{equation}
其中$\mathcal{F}$为回归树函数空间。第$m$步优化带正则化的损失函数：
\begin{equation}
\mathcal{L}^{(m)} = \sum_{t}\, l\!\left(C_t,\;\hat{C}_t^{(m-1)}+f_m(\mathbf{X}_t)\right) + \Omega(f_m),
\end{equation}
正则项$\Omega(f_m)=\gamma J + \frac{1}{2}\lambda\|\mathbf{w}\|^2$控制叶节点数$J$与叶权重$\mathbf{w}$的复杂度，抑制过拟合。XGBoost对特征交互具有较强捕捉能力，但在小样本时序中稳定性不及线性模型。

\item {\songti\bfseries{}集成模型（Ensemble）}：将AR2、二次趋势和多变量三类模型的预测值加权组合：
\begin{equation}
\hat{C}_t^{(\mathrm{ens})} = \omega_1 \hat{C}_t^{(\mathrm{AR2})} + \omega_2 \hat{C}_t^{(\mathrm{Trend})} + \omega_3 \hat{C}_t^{(\mathrm{Multi})},
\end{equation}
权重满足$\omega_1+\omega_2+\omega_3=1$，$\omega_j\geq 0$，通过验证窗口RMSE的倒数归一化确定：
\begin{equation}
\omega_j = \frac{1/\mathrm{RMSE}_j}{\sum_{j=1}^{3}1/\mathrm{RMSE}_j}.
\end{equation}
权重自动向误差更小的模型倾斜，在模型间存在互补信息时可降低整体预测方差。

\end{enumerate}

评价指标采用RMSE、MAE和MAPE，其定义为：
\begin{align}
\mathrm{RMSE} &= \sqrt{\frac{1}{n}\sum_{t=1}^{n}(C_t-\hat{C}_t)^2}, \\[4pt]
\mathrm{MAE}  &= \frac{1}{n}\sum_{t=1}^{n}|C_t-\hat{C}_t|, \\[4pt]
\mathrm{MAPE} &= \frac{1}{n}\sum_{t=1}^{n}\left|\frac{C_t-\hat{C}_t}{C_t}\right|\times 100\%.
\end{align}
RMSE对大误差惩罚更重，适合关注极端偏差的场景；MAE具有鲁棒性；MAPE提供无量纲的相对误差，便于跨样本比较。

为保证比较公平，本文在滚动一步预测中采用统一可比窗口：仅在各模型均可生成预测值的年份区间上计算误差指标。

\subsection{情景模拟与可解释性分析}
在预测比较基础上，本文设置基准、低碳和深度脱碳三类情景，对2023年后的路径进行模拟。情景以最后一期观测值$C_{T_0}$（$T_0=2022$年）为起点，按不同政策力度设定年度衰减路径：
\begin{equation}
C_t^{(s)} = C_{T_0} \cdot g^{(s)}(t - T_0), \quad t = T_0+1,\ldots,2035,
\end{equation}
其中$g^{(s)}(\cdot)$为情景$s$对应的路径函数。具体地：

\begin{itemize}
  \item {\songti\bfseries{}基准情景（Baseline）}：反映既有政策延续下的排放惯性，近期以较小速率下降，随后受技术积累影响降速略有加快：
  \begin{equation}
  g^{(\mathrm{base})}(k) = 1 - 0.0035k + 0.00015\cdot\max(k-3,0)^2.
  \end{equation}
  \item {\songti\bfseries{}低碳情景（Low-Carbon）}：反映能源结构优化与产业升级加速，排放以较快线性速率下降：
  \begin{equation}
  g^{(\mathrm{low})}(k) = 1 - 0.014k.
  \end{equation}
  \item {\songti\bfseries{}深度脱碳情景（Deep Decarbonization）}：反映能源替代与效率提升协同强化，以最快速率压减总量：
  \begin{equation}
  g^{(\mathrm{deep})}(k) = 1 - 0.024k.
  \end{equation}
\end{itemize}

情景参数锚定于国家能源转型目标导向与近年宏观运行区间，三类情景的GDP增速、煤炭占比年均变动和能源强度下降率设定详见正文表~\ref{tab:scenario_params}。

为增强预测结果的可解释性，本文基于XGBoost模型输出计算SHAP（SHapley Additive exPlanations）值。对于样本点$t$，特征$j$的SHAP值定义为：
\begin{equation}
\phi_j(t) = \sum_{S \subseteq \mathcal{F}\setminus\{j\}} \frac{|S|!\,(|\mathcal{F}|-|S|-1)!}{|\mathcal{F}|!}
\left[\hat{f}_{S\cup\{j\}}(\mathbf{x}_t^{S\cup\{j\}}) - \hat{f}_S(\mathbf{x}_t^S)\right],
\end{equation}
其中$\mathcal{F}$为全部特征集合，$S$为不含$j$的特征子集，$\hat{f}_S$为仅使用子集$S$中特征的模型预测。SHAP值满足效率性（所有特征贡献之和等于预测值与基线之差）、对称性与虚拟性等公理，是目前最具理论基础的模型解释框架之一。本文以各特征SHAP值的绝对值均值$\overline{|\phi_j|}$衡量全局重要性，实现”预测准确性”与”机制可解释性”的同框分析。

\subsection{方法适用性与局限说明}
本文方法适用于国家级中短期时间序列和总量路径识别场景。其优势在于结构清晰、变量含义明确、政策映射直接。局限在于：样本长度约束可能影响复杂模型表现；中介路径难以完成严格因果识别；情景模拟对参数设定仍有一定依赖。为降低这些局限的影响，本文采用滚动外推评价、跨模型比较和观测/预测边界区分等策略提高结论可信度。
